% =============================================================================
%  Compte-rendu de TP -- Cours 6 « DevOps pour le HPC » (M2 CHPS, CHPS0902)
%  Carnet de bord de la séance en pédagogie inversée (missions 1 à 5)
%
%  Compilation : pdflatex carnet-de-bord.tex (deux passes) ou latexmk
%  Les extraits de code et de logs sont lus directement dans le dépôt et dans
%  logs/ : rejouer une commande met le document à jour.
% =============================================================================
\documentclass[11pt,a4paper]{article}

% ------------------------------------------------------------- à renseigner
\newcommand{\etudiants}{Sylvain ROWE}
\newcommand{\depot}{\afaire{URL du dépôt}}
\newcommand{\gcpinfo}{\afaire{type de VM, partition}}
\newcommand{\spackinstalltime}{\afaire{durée}}
\newif\ifbrouillon \brouillontrue  % \brouillonfalse : masque la page « Reste à faire »
% Les scripts scripts/gitlab_evidence.py et scripts/gcp_evidence.sh écrivent leurs
% preuves dans logs/gitlab/ et logs/gcp/ : le document les inclut dès qu'elles existent.
\newif\ifgitlab \newif\ifgcp
% -----------------------------------------------------------------------------

\usepackage[T1]{fontenc}
\usepackage[utf8]{inputenc}
\usepackage{lmodern}
\usepackage{courier}
\usepackage[scaled=0.9]{helvet}
\usepackage[french]{babel}
\frenchsetup{SmallCapsFigTabCaptions=false, StandardItemLabels=true}
\usepackage[babel]{csquotes}
\addto\captionsfrench{\renewcommand{\tablename}{Tableau}}

\usepackage[a4paper,margin=2.5cm]{geometry}
\usepackage{microtype}
\usepackage{amsmath,amssymb}
\usepackage{graphicx}
\usepackage{booktabs,tabularx,array}
\usepackage[dvipsnames]{xcolor}
\usepackage{enumitem}
\usepackage{listings}
\usepackage{caption}
\usepackage{tikz}
\usetikzlibrary{positioning,arrows.meta,calc,fit,backgrounds,decorations.pathmorphing}
\usepackage[obeyspaces,spaces]{url}
\usepackage[hidelinks]{hyperref}

\captionsetup{font=small, labelfont=bf, labelsep=period, justification=justified,
  singlelinecheck=true, skip=6pt}
\captionsetup[table]{position=top}
\setlist{itemsep=0.1em, topsep=0.3em}
\setlength{\emergencystretch}{2.5em}

% ------------------------------------------------------------- code
\DeclareUrlCommand\code{\urlstyle{tt}}
\def\UrlBreaks{\do\/\do\-\do\_\do\.\do\=\do\,\do\:\do\;}
\newcommand{\fic}[1]{\code{#1}}
\newcolumntype{L}{>{\raggedright\arraybackslash}X}

\definecolor{gristrait}{gray}{0.6}
\definecolor{gristexte}{gray}{0.45}
\lstset{
  basicstyle=\ttfamily\scriptsize, columns=fullflexible, keepspaces=true,
  breaklines=true, breakatwhitespace=false, showstringspaces=false, upquote=true,
  frame=lines, rulecolor=\color{gristrait}, framesep=3pt, xleftmargin=0pt,
  aboveskip=8pt, belowskip=8pt, captionpos=b, abovecaptionskip=6pt,
  tabsize=4, extendedchars=true,
  postbreak=\mbox{\textcolor{gristexte}{$\hookrightarrow$}\space},
  commentstyle=\color{gristexte}, keywordstyle=\bfseries, stringstyle={},
  literate={é}{{\'e}}1 {è}{{\`e}}1 {ê}{{\^e}}1 {ë}{{\"e}}1 {à}{{\`a}}1 {â}{{\^a}}1
           {ç}{{\c c}}1 {ô}{{\^o}}1 {î}{{\^i}}1 {ï}{{\"i}}1 {û}{{\^u}}1 {ù}{{\`u}}1
           {É}{{\'E}}1 {È}{{\`E}}1 {À}{{\`A}}1 {Ç}{{\c C}}1 {œ}{{\oe}}1 {Œ}{{\OE}}1
           {«}{{\guillemotleft}}1 {»}{{\guillemotright}}1 {×}{{$\times$}}1
           {→}{{$\to$}}1 {’}{{'}}1 {…}{{...}}3
}
\lstdefinestyle{numerote}{numbers=left, numberstyle=\tiny\color{gristexte}, numbersep=6pt}
\lstdefinelanguage{yaml}{morecomment=[l]{\#}, sensitive=true,
  morekeywords={workflow,stages,variables,default,image,interruptible,before_script,
    script,after_script,rules,if,when,needs,job,artifacts,paths,expire_in,reports,junit,
    stage,extends,tags,timeout,retry,max,exit_codes,allow_failure,spack,specs,view,
    concretizer,unify,packages,all,target,externals,spec,prefix,buildable}}
\renewcommand{\lstlistingname}{Listing}

% ------------------------------------------------------------- éléments à compléter
\newcommand{\afaire}[1]{\textcolor{red!70!black}{[À compléter : #1]}}
% emplacement d'une capture : remplacer par \includegraphics[width=\linewidth]{...}
\newcommand{\capture}[2]{%
  \begin{tikzpicture}
    \draw[red!70!black, dashed, rounded corners=2pt] (0,0) rectangle (\linewidth,#1);
    \node[text width=0.85\linewidth, align=center, red!70!black, font=\small]
      at (0.5\linewidth,0.5*#1) {Capture à insérer\\[2pt] #2};
  \end{tikzpicture}}

% ------------------------------------------------------------- figures (palette Okabe-Ito)
\definecolor{oiBleu}{RGB}{0,114,178}
\definecolor{oiCiel}{RGB}{86,180,233}
\definecolor{oiOrange}{RGB}{230,159,0}
\definecolor{oiVert}{RGB}{0,158,115}
\definecolor{oiVermillon}{RGB}{213,94,0}
\definecolor{oiPourpre}{RGB}{204,121,167}
\tikzset{
  job/.style={draw=#1, fill=#1!18, rounded corners=2pt, minimum width=2.75cm,
    minimum height=0.62cm, font=\ttfamily\footnotesize, inner sep=2pt},
  quand/.style={font=\sffamily\scriptsize, text=black!70},
  fleche/.style={-{Stealth[length=2.2mm]}, thick},
  boite/.style={draw=black!55, rounded corners=3pt, inner sep=6pt},
}
\newcommand{\jobtxt}[2]{{\ttfamily\scriptsize #1}\\[-1pt]{\sffamily\tiny\color{black!70}#2}}

\IfFileExists{logs/gitlab/meta.tex}{\gitlabtrue\input{logs/gitlab/meta.tex}}{}
\IfFileExists{logs/gcp/meta.tex}{\gcptrue\input{logs/gcp/meta.tex}}{}
\ifgitlab\ifgcp\brouillonfalse\fi\fi

% =============================================================================
\begin{document}

% ------------------------------------------------------------- page de titre
\begin{titlepage}
\centering
\vspace*{0.4cm}
{\large Université de Reims Champagne-Ardenne\par}
\vspace{0.15cm}
{\large Master 2 Calcul Haute Performance et Simulation (CHPS)\par}
\vspace{4.2cm}
{\Large\bfseries Compte-rendu de TP\par}
\vspace{0.4cm}
{\large Cours 6 : DevOps pour le HPC, séance en pédagogie inversée\par}
\vspace{1.8cm}
{\LARGE\bfseries Du dépôt Git au cluster :\par\vspace{0.25cm}
intégration continue, Slurm et Spack\par}
\vspace{0.8cm}
{\large carnet de bord des missions 1 à 5,\par
projet \texttt{mpiflow}, tag \texttt{seance-inversee}\par}
\vfill
{\large
\begin{tabular}{r@{\quad}l}
Étudiant(s) : & \etudiants\\
Formation : & M2 CHPS\\
UE : & CHPS0902\\
Enseignant : & Luiz Angelo Steffenel\\
Date : & \today\\
\end{tabular}\par}
\vspace{2.2cm}
\end{titlepage}

% ------------------------------------------------------------- reste à faire
\ifbrouillon
\thispagestyle{empty}
\section*{\textcolor{red!70!black}{Reste à faire avant le rendu}}
{\color{red!70!black}\small Cette page disparaît d'elle-même quand les deux scripts ci-dessous
ont produit leurs preuves ; les mentions \afaire{...} sont alors remplies.}
\begin{enumerate}
\item Sur votre machine, depuis le dépôt (branche \code{main}) :
\code{GITLAB_TOKEN=glpat-... python3 scripts/gitlab_evidence.py}
\ifgitlab\enskip\textcolor{OliveGreen}{fait}\fi
\item Sur le nœud de connexion du cluster GCP, depuis le dépôt cloné dans \fic{/home} :
\code{bash scripts/gcp_evidence.sh} \ifgcp\enskip\textcolor{OliveGreen}{fait}\fi
\item Recompiler (\code{latexmk -g} et \code{python3 scripts/tex2md.py}, faits par les scripts
si LaTeX et pandoc sont installés), commiter \fic{logs/}, \fic{spack/}, le PDF et le Markdown, puis
\code{git tag -f seance-inversee && git push -f origin seance-inversee}.
\item \code{terraform destroy} et vérification de la console GCP.
\end{enumerate}
\clearpage
\setcounter{page}{1}
\fi

% =============================================================================
\section{Introduction}
% =============================================================================

Un code de calcul ne vit pas seulement sur le poste du développeur : il est compilé, testé,
empaqueté puis exécuté sur un cluster partagé, où tout passe par l'ordonnanceur Slurm et où
l'environnement logiciel doit être identique sur chaque nœud. Le DevOps appliqué au HPC
automatise cette chaîne, mais ses outils ne se transposent pas tels quels depuis le cloud :
\enquote{CI $\neq$ exécution cluster}.

Cette séance en pédagogie inversée consistait à le vérifier par l'expérience, en quatre
missions : faire exécuter par une autre machine la compilation et le test d'un code MPI
(mission~1), conditionner les jobs à des événements et protéger les secrets (mission~2), faire
franchir au pipeline la frontière entre le runner CI et Slurm (mission~3), et versionner
l'environnement logiciel avec Spack (mission~4). Ce carnet de bord constitue la mission~5 : une
section par mission, avec les réponses aux questions de cadrage, les preuves d'exécution, les
erreurs rencontrées et un bilan, puis un auto-diagnostic.

Pour chaque question, nous avons cherché une réponse fondée sur une expérience plutôt que sur
les diapositives : plusieurs réponses en sont sorties nuancées, voire contraires à notre
intuition de départ.

% =============================================================================
\section{Matériels et méthodes}
% =============================================================================

\subsection{Environnement}

Pour pouvoir casser librement les choses (faire échouer un test, saturer une file, simuler la
panne d'un nœud), nous avons d'abord reproduit chaque mission sur un poste Linux, avec un
petit cluster Slurm complet et un émulateur de GitLab CI (tableau~\ref{tab:env}). Le
mini-cluster comprend le contrôleur, la base de comptabilité (sans laquelle \code{sacct} ne
fonctionne pas) et deux nœuds de 2 CPU, comme deux VM \code{e2-standard-2}. Les mêmes fichiers
sont ensuite exécutés sur GitLab.com et sur le cluster GCP déployé par Terraform au Cours~5.

\begin{table}[ht]
\centering\small
\caption{Environnements d'exécution.}\label{tab:env}
\begin{tabularx}{\linewidth}{@{}>{\raggedright\arraybackslash}p{3.2cm} L >{\raggedright\arraybackslash}p{2.6cm}@{}}
\toprule
Composant & Version et configuration & Missions\\ \midrule
Poste de reproduction & Ubuntu 24.04, 2 vCPU, 8 Go ; GCC 13.3, CMake 3.28, Open MPI 4.1.6 & toutes\\
Analyse statique & cppcheck 2.13, shellcheck 0.9, actionlint 1.7.12 & 2, 3\\
Émulateur GitLab CI & \code{gitlab-ci-local} 4.76.1 : validation par le schéma officiel, évaluation des \code{rules}, exécution des jobs & 1, 2\\
Mini-cluster Slurm & Slurm 23.11.4 : \code{slurmctld}, \code{slurmdbd} et MariaDB, deux \code{slurmd} (\code{node1}, \code{node2}, 2 CPU), partition \code{debug}, utilisateur sans privilèges & 3\\
Spack & 1.2.2 (tag de version), dépôt de paquets \code{builtin}, solveur clingo 5.8 & 4\\
GitLab.com & runners partagés, projet \depot & 1 à 3\\
Google Cloud & cluster Slurm du Cours 5 : \gcpinfo & 3, 4\\
\bottomrule
\end{tabularx}
\end{table}

\subsection{Le dépôt \texttt{mpiflow}}

Le dépôt de travail (tableau~\ref{tab:depot}) contient un code de test MPI, son build CMake et
tous les fichiers produits pendant les missions. Le pipeline final est résumé par la
figure~\ref{fig:pipeline} ; le fichier complet est donné en annexe~A.

\begin{table}[ht]
\centering\small
\caption{Contenu du dépôt.}\label{tab:depot}
\begin{tabularx}{\linewidth}{@{}>{\ttfamily\footnotesize\raggedright\arraybackslash}p{5.5cm} L c@{}}
\toprule
\normalfont Fichier & Rôle & Mission\\ \midrule
tests/test\_mpi.c & test MPI : Bcast, Allreduce, taille du communicateur & 1\\
CMakeLists.txt, tests/CMakeLists.txt & build CMake, tests CTest à 1 et 4 rangs & 1\\
.gitlab-ci.yml & pipeline : build, test, lint, package, secrets, Slurm & 1 à 3\\
.github/workflows/ci.yml & équivalent GitHub Actions & 1, 2\\
scripts/fake\_docker\_login.sh & simulation de \texttt{docker login -{}-password-stdin} & 2\\
slurm/job\_ci.sh & job Slurm : 2 nœuds, MPI, \texttt{SUCCESS} & 3\\
slurm/submit\_and\_poll.sh & soumission, suivi par \texttt{sacct}, verdict & 3\\
slurm/ci\_remote.sh & variante SSH : build et soumission à distance & 3\\
spack/spack.yaml, spack/spack.lock & environnement Spack versionné & 4\\
logs/ & sorties brutes citées dans ce document & --\\
\bottomrule
\end{tabularx}
\end{table}

\begin{figure}[ht]
\centering
\begin{tikzpicture}[x=1cm,y=1cm,
    jb/.style={draw=#1, fill=#1!18, rounded corners=2pt, minimum width=2.6cm,
      minimum height=0.95cm, align=center, inner sep=2pt}]
  \foreach \x/\s in {0/lint,3.3/build,6.6/test,9.9/package,13.2/test\_hpc}
    \node[font=\sffamily\small\bfseries, text=black!65] at (\x,1.6) {\s};
  \draw[black!25] (-1.4,1.3) -- (14.6,1.3);
  \node[jb=oiPourpre] (lint) at (0,0) {\jobtxt{lint}{MR uniquement}};
  \node[jb=oiCiel] (build) at (3.3,0) {\jobtxt{build}{tous les pipelines}};
  \node[jb=oiCiel] (test) at (6.6,0) {\jobtxt{test}{tous les pipelines}};
  \node[jb=oiOrange] (pkg) at (9.9,0) {\jobtxt{package}{tags vX.Y.Z}};
  \node[jb=oiCiel] (reg) at (9.9,-1.5) {\jobtxt{registry\_login}{tous ; secret protégé}};
  \node[jb=oiVert] (hpc) at (13.2,0) {\jobtxt{test\_slurm}{nuit, ou manuel sur main}};
  \node[jb=oiVert] (wait) at (13.2,-1.5) {\jobtxt{test\_slurm\_wait}{nuit, ou manuel sur main}};
  \node[jb=oiVert, dashed] (ssh) at (13.2,-3.0) {\jobtxt{test\_slurm\_ssh}{si clé SSH déclarée}};
  % dépendances needs
  \draw[fleche] (build) -- (test);
  \draw[fleche] (build.north) to[out=25,in=155]
    node[below, quand, pos=0.5] {artefact \texttt{build/}} (pkg.north);
  \draw[fleche, dashed] (test) -- (pkg);
  \draw[thick, dashed] (test.south) -- (6.6,-0.78) -- (11.6,-0.78);
  \draw[fleche, dashed] (11.6,-0.78) |- (hpc.west);
  \draw[fleche, dashed] (11.6,-0.78) |- (wait.west);
  % légende
  \begin{scope}[shift={(-1.3,-4.1)}, font=\sffamily\scriptsize]
    \draw[fleche] (0,0) -- (0.7,0) node[right] {\texttt{needs} avec artefacts};
    \draw[fleche,dashed] (3.7,0) -- (4.4,0) node[right] {\texttt{needs} sans artefact};
    \node[jb=oiCiel, minimum width=0.5cm, minimum height=0.3cm] at (7.7,0) {};
    \node[right] at (8.0,0) {CI rapide};
    \node[jb=oiOrange, minimum width=0.5cm, minimum height=0.3cm] at (9.7,0) {};
    \node[right] at (10.0,0) {version};
    \node[jb=oiVert, minimum width=0.5cm, minimum height=0.3cm] at (11.6,0) {};
    \node[right] at (11.9,0) {cluster (Slurm)};
    \node[jb=oiPourpre, minimum width=0.5cm, minimum height=0.3cm] at (14.2,0) {};
    \node[right] at (14.5,0) {MR};
  \end{scope}
\end{tikzpicture}
\caption{Pipeline final. Chaque colonne est un stage ; les flèches représentent les
dépendances \texttt{needs}, sous chaque job figure la condition qui le crée. \texttt{lint},
\texttt{registry\_login} et \texttt{test\_slurm\_ssh} ont \texttt{needs: []} et démarrent dès la
création du pipeline. Les jobs Slurm attendent \texttt{test} sans récupérer son binaire, compilé
pour l'image Debian et non pour le cluster.}
\label{fig:pipeline}
\end{figure}

\subsection{Démarche et preuves}

Chaque commande citée a été exécutée à travers un petit enregistreur qui écrit la commande, sa
sortie complète et son code retour dans un fichier de \fic{logs/}. Les listings de ce document
incluent ces fichiers tels quels ; les rares coupes sont signalées par \code{[...]}. Les preuves
produites sur GitLab.com et sur le cluster GCP sont collectées de la même façon, par
\fic{scripts/gitlab_evidence.py} (API de GitLab) et \fic{scripts/gcp_evidence.sh} (sur le nœud
de connexion), dans \fic{logs/gitlab/} et \fic{logs/gcp/}.

% =============================================================================
\section{Mission 1 : un pipeline minimal}
% =============================================================================

\subsection{Réalisation}

Le test \fic{tests/test_mpi.c} (annexe~C) vérifie trois choses : la valeur 42 diffusée par
\code{MPI_Bcast} arrive sur tous les rangs, tous les rangs participent (\code{MPI_Allreduce} de la
somme des rangs) et, si la variable \code{EXPECTED_RANKS} est définie, \code{MPI_COMM_WORLD} a bien
cette taille. Le rang 0 affiche aussi le nœud de chaque rang, ce qui servira de preuve
d'exécution multi-nœuds en mission~3. Les vérifications sont explicites et renvoient un code
d'erreur, au lieu d'utiliser \code{assert} comme le test des diapositives : nous avons constaté
que cette assertion disparaît en compilation \code{Release} (section~\ref{sec:m1-erreurs}).

CTest déclare deux tests, à 1 et 4 rangs, à partir des variables \code{MPIEXEC_*} de
\code{FindMPI}. Le pipeline (listing~\ref{lst:m1-yaml}) enchaîne deux jobs dans l'image
\code{gcc:13-bookworm} : \code{build} publie \fic{build/} comme artefact, \code{test} le récupère via
\code{needs: [build]}, lance \code{ctest} (le rapport JUnit s'affiche dans GitLab) puis
\code{mpirun -np 4}.

\lstinputlisting[language=yaml, style=numerote, firstnumber=47, linerange={47-70},
  caption={Jobs \texttt{build} et \texttt{test} (\texttt{.gitlab-ci.yml}, lignes 47 à 70).},
  label=lst:m1-yaml]{.gitlab-ci.yml}

\subsection{Preuves d'exécution}

Lancé en root sur une machine à 2 cœurs, \code{mpirun} refuse deux fois de démarrer avant de
réussir avec les deux options demandées par le sujet (listing~\ref{lst:m1-mpirun}). CTest, configuré
avec ces options, passe ses deux tests (listing~\ref{lst:m1-ctest}). Le même pipeline, exécuté
par l'émulateur avec des artefacts isolés entre jobs, donne \code{PASS build} et \code{PASS test}
(variante A du tableau~\ref{tab:needs}).

\lstinputlisting[caption={\texttt{mpirun} en root sur 2 cœurs : refus, refus, puis succès (\texttt{logs/extraits/m1-mpirun.txt}).},
  label=lst:m1-mpirun]{logs/extraits/m1-mpirun.txt}
\lstinputlisting[caption={CTest avec les options du CI (\texttt{logs/m1-ctest.txt}).},
  label=lst:m1-ctest]{logs/m1-ctest.txt}

\ifgitlab
Sur GitLab.com, le même pipeline s'exécute sur les runners partagés (listing~\ref{lst:gl-m1}).
\lstinputlisting[caption={Pipeline de branche sur GitLab.com : jobs et fin du log du job \texttt{test} (\texttt{logs/gitlab/m1-pipeline-main.txt}).},
  label=lst:gl-m1]{logs/gitlab/m1-pipeline-main.txt}
\else
\afaire{exécution de \texttt{scripts/gitlab\_evidence.py} : le log du pipeline de branche sur
GitLab.com s'insérera ici}
\fi

\subsection{Questions de cadrage}

\subsubsection{Deux jobs d'un même stage échouent : le pipeline continue-t-il ?}

Nous avons construit un pipeline de trois builds parallèles, dont deux échouent, suivis d'un
stage de test (listing~\ref{lst:m1-stage}). On voit que \code{build_intel} échoue en premier sans
interrompre \code{build_gcc}, qui échoue 2\,s plus tard : les jobs d'un même stage sont
indépendants et vont tous au bout. On voit aussi que le job \code{test} n'apparaît pas : il n'a
jamais été lancé, tandis que \code{cleanup}, déclaré \code{when: on_failure}, s'est exécuté. La
documentation le confirme : les jobs d'un stage ne démarrent qu'après la \emph{réussite} des jobs
du stage précédent. Il existe trois exceptions : \code{allow_failure: true} rend un échec non
bloquant, les jobs \code{when: on_failure} ou \code{always} servent au nettoyage et aux
notifications, et un job relié par \code{needs:} à des jobs qui ont réussi ne dépend pas des
autres. À l'inverse, \code{workflow:auto_cancel:on_job_failure: all} annule tout au premier échec.

\lstinputlisting[firstline=26, lastline=38,
  caption={Deux échecs dans le stage \texttt{build} (\texttt{logs/m1-stage-fail.txt}).},
  label=lst:m1-stage]{logs/m1-stage-fail.txt}

Ce comportement est souhaitable pour trois raisons. Il fait échouer tôt : tester ou empaqueter
un binaire qui n'a pas compilé n'a pas de sens, et en HPC l'étape suivante coûte des heures de
file d'attente et de quota. Il empêche de publier un artefact non validé : c'est aussi pourquoi
notre job \code{package} déclare \code{needs: [build, test]}, car \code{needs: [build]} seul
contournerait l'ordre des stages et publierait même si les tests échouent. Enfin, laisser les
jobs parallèles aller au bout donne toutes les causes d'échec (ici gcc et intel) dans un seul
pipeline.

\subsubsection{Quelle différence entre mpirun dans le conteneur CI et une exécution sur cluster ?}

\begin{table}[ht]
\centering\small
\caption{Exécution MPI dans le conteneur CI et sur le cluster.}\label{tab:ci-cluster}
\begin{tabularx}{\linewidth}{@{}>{\raggedright\arraybackslash}p{2.6cm} L L@{}}
\toprule
 & Conteneur CI (runner partagé) & Cluster (Slurm)\\ \midrule
Placement & 1 machine de 2 à 4 vCPU : 4 rangs sur 1 hôte, en surréservation & rangs répartis sur plusieurs nœuds alloués\\
Lancement & \code{mpirun} lance lui-même les processus, sans allocation & \code{srun} dans une allocation Slurm, via PMIx\\
Communications & mémoire partagée uniquement & réseau entre nœuds, mémoire partagée dans un nœud\\
Environnement & Open MPI 4.1 de Debian, utilisateur root & MPI du cluster (Spack ou modules), utilisateur sans privilèges\\
Ce que ça valide & la logique du code sur peu de rangs, à chaque push & l'intégration du code, du MPI, du lanceur et du réseau ; les performances\\
\bottomrule
\end{tabularx}
\end{table}

Le tableau~\ref{tab:ci-cluster} résume les différences ; nos expériences les rendent concrètes.
En CI, il a fallu \code{--allow-run-as-root} (le conteneur tourne en root) et
\code{--oversubscribe} (4 rangs pour 2 cœurs), et le test affiche \code{4 ranks on 1 host(s)} :
aucune communication n'a traversé un réseau et les temps mesurés n'ont aucun sens. Sur le
mini-cluster, \emph{le même binaire} a échoué de deux façons invisibles en CI : lancé par
\code{srun} sans PMIx, Open MPI refuse de démarrer, et réparti sur deux nœuds, un rang a planté
pendant que les autres restaient bloqués jusqu'au wall-time (section~\ref{sec:m3-preuves}). Enfin,
le binaire produit dans le conteneur est lié à l'Open MPI de l'image Debian, pas à celui du
cluster : les jobs HPC recompilent sur le nœud de connexion. On en déduit que le CI joue le rôle
d'un test unitaire du code MPI, et le cluster celui d'un test d'intégration du couple code et
environnement.

\subsubsection{Pourquoi le test échoue-t-il si l'on oublie needs: alors que le build a réussi ?}

L'expérience nous a obligés à nuancer la question : supprimer \code{needs:} \emph{ne fait pas}
échouer le test sur GitLab (tableau~\ref{tab:needs}, variante B). Chaque job démarre dans un
environnement neuf (nouveau conteneur, nouveau clone) où \fic{build/} n'existe pas, mais par
défaut GitLab télécharge les artefacts de tous les jobs des stages précédents. Le test n'échoue
que s'il ne reçoit pas l'artefact : avec \code{needs: []}, qui le fait démarrer sans attendre
(variante C, listing~\ref{lst:m1-needs}), ou s'il est placé dans le même stage que le build
(variante D). L'échec \enquote{needs oublié} est en revanche systématique sur GitHub Actions, où
des jobs sans \code{needs:} tournent en parallèle et ne partagent aucun fichier.

\begin{table}[ht]
\centering\small
\caption{Effet de \texttt{needs} sur le job \texttt{test} (\texttt{logs/m1-needs.txt}).}\label{tab:needs}
\begin{tabularx}{\linewidth}{@{}l L c@{}}
\toprule
Variante du job \texttt{test} & Comportement observé & Résultat\\ \midrule
A. \code{needs: [build]} & artefact importé, \code{MPI test passed} & réussi\\
B. \code{needs:} supprimé & artefact importé quand même (stage précédent) & réussi\\
C. \code{needs: []} & démarre sans attendre : \code{cd: build: No such file or directory} & échec\\
D. \code{needs:} supprimé, même stage que \code{build} & jobs parallèles, pas d'artefact : même erreur & échec\\
\bottomrule
\end{tabularx}
\end{table}

\lstinputlisting[firstline=27, lastline=36,
  caption={Variante C : \texttt{needs: []} (\texttt{logs/m1-needs.txt}).},
  label=lst:m1-needs]{logs/m1-needs.txt}

Sur GitLab, \code{needs:} sert donc à autre chose : le test démarre dès la fin de \code{build} sans
attendre le reste du stage, et seuls les artefacts listés sont téléchargés (sinon les
\fic{build/} de plusieurs compilateurs arriveraient tous au même chemin). L'expérience a révélé un
second piège : \code{ctest} indique \code{Test project .../builds/build/build}, c'est-à-dire le
répertoire du job \code{build}. Un arbre de build CMake contient des chemins absolus et n'est pas
relocalisable ; cela fonctionne sur GitLab.com parce que tous les jobs tournent au même chemin
(\code{/builds/<groupe>/<projet>}), mais nous l'avons payé avec \code{cpack} en mission~2.

\subsection{Erreurs rencontrées et résolution}\label{sec:m1-erreurs}

\begin{table}[ht]
\centering\small
\caption{Erreurs de la mission 1.}\label{tab:m1-erreurs}
\begin{tabularx}{\linewidth}{@{}>{\raggedright\arraybackslash}p{4.2cm} L L@{}}
\toprule
Symptôme & Cause & Résolution\\ \midrule
\code{mpirun has detected an attempt to run as root} & le conteneur CI s'exécute en root & \code{--allow-run-as-root}\\
\code{not enough slots available} & 4 rangs demandés, 2 cœurs & \code{--oversubscribe}, passé à CTest par \code{MPIEXEC_PREFLAGS}\\
Le test des diapositives passe malgré un bug volontaire (listing~\ref{lst:m1-ndebug}) & \code{Release} ajoute \code{-DNDEBUG}, qui supprime \code{assert()} & vérifications explicites et code retour ; vérifier qu'un test \emph{peut} échouer\\
Jobs en échec \code{FAIL 100} dans l'émulateur & deux \code{apt-get} parallèles sur la même machine (exécuteur \emph{shell}) & dépendances préinstallées ; résultats écartés\\
Non observé (pas de Docker en local) & Open MPI en conteneur Docker : copie mémoire CMA impossible & \code{OMPI_MCA_btl_vader_single_copy_mechanism=none}, posé préventivement\\
\bottomrule
\end{tabularx}
\end{table}

\lstinputlisting[caption={Le test des diapositives avec un bug volontaire : détecté sans
\texttt{-DNDEBUG}, invisible en \texttt{Release} (\texttt{logs/m1-ndebug.txt}).},
  label=lst:m1-ndebug]{logs/m1-ndebug.txt}

\subsection{Bilan}
\begin{description}[font=\normalfont\bfseries, leftmargin=0pt, itemsep=0.2em]
\item[Ce qui fonctionne :] build CMake, CTest à 1 et 4 rangs, pipeline validé par le schéma
GitLab et exécuté de bout en bout par l'émulateur.
\item[Ce qui ne fonctionne pas :] \ifgitlab rien de bloquant ; le pipeline ne teste toutefois
que 4 rangs sur une seule machine. \else le comportement des runners de GitLab.com reste à relever
(pas de Docker sur le poste de reproduction).\fi
\item[À refaire autrement :] écrire d'abord un test qui échoue avant de lui faire confiance ;
c'est ce qui a révélé le piège de \code{NDEBUG}.
\end{description}

% =============================================================================
\section{Mission 2 : déclencheurs, règles et secrets}
% =============================================================================

\subsection{Réalisation}

Trois jobs s'ajoutent au pipeline (figure~\ref{fig:pipeline}). Le job \code{lint} lance
\code{cppcheck} avec \code{--error-exitcode=1}, uniquement dans les pipelines de merge request. Le
job \code{package} ne s'exécute que sur les tags \code{vX.Y.Z} : il archive le binaire, le README et
un fichier \fic{BUILDINFO.txt} (projet, tag, commit, URL du pipeline, date), puis calcule le
SHA-256 de l'archive. Le job \code{registry_login} lit la variable masquée et protégée
\code{FAKE_REGISTRY_TOKEN} et la transmet par l'entrée standard à
\fic{scripts/fake_docker_login.sh}, qui reproduit le contrat de \code{docker login --password-stdin}
sans démon Docker. Un bloc \code{workflow:rules} choisit le type de pipeline
(listing~\ref{lst:m2-workflow}) et les jobs sans condition reçoivent \code{rules: - when: on_success}.

\lstinputlisting[language=yaml, style=numerote, firstnumber=15, linerange={15-23},
  caption={Choix du type de pipeline (\texttt{.gitlab-ci.yml}, lignes 15 à 23).},
  label=lst:m2-workflow]{.gitlab-ci.yml}

\subsection{Preuves d'exécution}

L'émulateur évalue les règles pour chaque déclencheur (tableau~\ref{tab:rules}) : \code{lint}
n'existe que dans le pipeline de MR, \code{package} que pour un tag \code{vX.Y.Z} (pas pour
\code{seance-inversee}, listing~\ref{lst:m2-regex}), et les jobs Slurm que si un runner HPC est
déclaré. Sur une branche contenant une erreur volontaire, \code{cppcheck} échoue avec le code 1
(listing~\ref{lst:m2-lint}). Le pipeline du tag \code{v0.1.0} produit une archive dont le contenu
et la somme de contrôle se vérifient (listing~\ref{lst:m2-package}).

\begin{table}[ht]
\centering\small
\caption{Jobs créés selon le déclencheur (\texttt{logs/m2-rules.txt}).}\label{tab:rules}
\begin{tabularx}{\linewidth}{@{}L c c c c c@{}}
\toprule
Pipeline & \texttt{lint} & \texttt{build}, \texttt{test} & \texttt{package} & \texttt{registry\_login} & \texttt{test\_slurm*}\\ \midrule
push sur une branche & -- & oui & -- & oui & --\\
merge request & oui & oui & -- & oui & --\\
tag \code{v1.2.0} & -- & oui & oui & oui & --\\
tag \code{seance-inversee} & -- & oui & -- & oui & --\\
planifié sur \code{main}, \code{HPC_ENABLED=true} & -- & oui & -- & oui & oui\\
push sur \code{main}, \code{HPC_ENABLED=true} & -- & oui & -- & oui & manuel\\
\bottomrule
\end{tabularx}
\end{table}

\lstinputlisting[caption={Quels tags déclenchent \texttt{package} ? (\texttt{logs/m2-regex.txt}).},
  label=lst:m2-regex]{logs/m2-regex.txt}
\lstinputlisting[firstline=4, caption={Lint sur une branche contenant une erreur volontaire (\texttt{logs/m2-lint.txt}).},
  label=lst:m2-lint]{logs/m2-lint.txt}
\lstinputlisting[caption={Pipeline du tag \texttt{v0.1.0} : archive, contenu et somme de contrôle (\texttt{logs/extraits/m2-package.txt}).},
  label=lst:m2-package]{logs/extraits/m2-package.txt}

Pour les secrets, le listing~\ref{lst:m2-secret} montre qu'un utilisateur sans privilèges lit
en clair, avec \code{ps}, le mot de passe passé en argument à un processus de root, alors que
\code{--password-stdin} ne laisse rien dans la ligne de commande. L'émulateur n'applique ni le
masquage ni la protection des variables : ces deux mécanismes ne se vérifient que sur GitLab.com
(figure~\ref{fig:cap-m2}).

\lstinputlisting[caption={Mot de passe en argument (\texttt{-p}) contre entrée standard (\texttt{logs/m2-secret.txt}).},
  label=lst:m2-secret]{logs/m2-secret.txt}

\ifgitlab
Sur GitLab.com (listings~\ref{lst:gl-reg} à~\ref{lst:gl-tag}), la variable apparaît
\code{[MASKED]} sur \code{main}, branche protégée, et vide dans le pipeline de la merge request ;
le commit contenant l'erreur volontaire ne crée qu'un pipeline, de type merge request, où
\code{lint} échoue ; le tag \code{v0.1.0} produit l'archive.
\lstinputlisting[caption={Variable masquée sur \texttt{main} (\texttt{logs/gitlab/m2-registry-main.txt}).},
  label=lst:gl-reg]{logs/gitlab/m2-registry-main.txt}
\lstinputlisting[caption={Merge request avec erreur volontaire : pipelines créés, lint, variable protégée absente (\texttt{logs/gitlab/m2-mr.txt}).},
  label=lst:gl-mr]{logs/gitlab/m2-mr.txt}
\lstinputlisting[caption={Pipeline du tag \texttt{v0.1.0} sur GitLab.com (\texttt{logs/gitlab/m2-tag.txt}).},
  label=lst:gl-tag]{logs/gitlab/m2-tag.txt}
\else
\afaire{exécution de \texttt{scripts/gitlab\_evidence.py} : masquage sur \texttt{main}, merge
request avec lint en échec et pipeline du tag \texttt{v0.1.0} s'inséreront ici}
\fi

\subsection{Questions de cadrage}

\subsubsection{rules: ou only:/except: : lequel écrire, et où retrouver l'autre ?}

Il faut écrire \code{rules:}. La documentation de GitLab classe désormais \code{only} et
\code{except} parmi les mots-clés dépréciés et renvoie vers \code{rules}. On les retrouve dans les
dépôts et gabarits anciens, comme le pipeline \enquote{hérité} des diapositives : il faut savoir
les lire et les traduire (\code{only: merge_requests} devient
\code{if: $CI_PIPELINE_SOURCE == "merge_request_event"}). Les deux ne se combinent pas dans un même
job : le validateur de GitLab rejette la configuration (\code{key may not be used with rules}),
alors que l'émulateur, plus laxiste, l'accepte.

\code{rules:} est plus expressif : conditions sur des variables, \code{changes}, \code{exists},
première règle vraie gagnante, \code{when: manual} ou \code{never} et \code{allow_failure} par règle ;
nous nous en servons pour rendre les jobs HPC manuels sur \code{main} et automatiques la nuit.
Mélanger les deux styles reste un piège, car un job sans règle vaut implicitement
\code{except: merge_requests}. Avec un \code{lint} limité aux MR et des jobs \code{build} et \code{test}
sans règle, chaque push sur une branche ayant une MR ouverte crée deux pipelines : un pipeline
de branche (build et test) et un pipeline de MR (lint seul). D'où le \code{workflow:rules} du
listing~\ref{lst:m2-workflow}.

\subsubsection{Que devient une variable Protected sur une branche non protégée ?}

Elle n'est pas injectée : le job démarre avec une variable vide, sans erreur. Notre job affiche
alors \code{FAKE_REGISTRY_TOKEN vide}, et \code{fake_docker_login.sh} refuse un secret vide (code
retour 1) ; un vrai job de déploiement doit échouer ainsi, explicitement, plutôt que de continuer
sans identifiant.

Le risque couvert est l'exfiltration par modification du pipeline. Toute personne qui peut
pousser une branche peut aussi modifier le \fic{.gitlab-ci.yml} de cette branche : un simple
\code{echo $TOKEN | base64} ou un \code{curl} vers un serveur externe contourne le masquage, qui ne
fait que remplacer la valeur exacte dans le log (la documentation précise d'ailleurs que le
masquage n'est pas une garantie). La protection réserve les secrets aux branches et tags
protégés, où seuls les mainteneurs poussent ou fusionnent, donc à du code relu en merge request.
Pour notre pipeline, la clé SSH du cluster et le futur jeton de buildcache sont donc protégés, et
le runner \code{hpc-login} doit lui aussi être déclaré protégé : sinon, n'importe quelle branche
pourrait exécuter du code arbitraire sur le nœud de connexion.

\subsubsection{Pourquoi ne jamais écrire un secret dans le .gitlab-ci.yml, même temporairement ?}

Parce que \enquote{temporairement} n'existe pas dans Git. Un secret commité reste dans
l'historique même après le commit qui le supprime, ainsi que dans tous les clones, forks, diffs de
MR et caches de GitLab ; l'effacer exige de réécrire l'historique, et le secret doit de toute
façon être considéré comme compromis et révoqué. Le fichier est en outre lisible par quiconque
a accès au dépôt (tout le monde si le projet est public, et des robots parcourent les forges à la
recherche de secrets). Enfin, une valeur écrite en clair n'est pas masquée : le masquage ne
s'applique qu'aux variables déclarées dans \emph{Settings > CI/CD}, et GitLab recopie chaque ligne
de script dans le log du job. Passée en argument (\code{-p}), elle apparaît en plus dans \code{ps}
pour les autres utilisateurs de la machine (listing~\ref{lst:m2-secret}).

Les variables de l'interface ont leurs propres contraintes : une valeur masquée doit tenir sur une
ligne, sans espace, et compter au moins 8 caractères (une fausse valeur trop courte comme
\code{1234} ne peut pas être masquée). Une variable de type \emph{File} contient le chemin d'un
fichier temporaire, ce qui convient à une clé SSH que \code{env} n'affichera pas.

\subsection{Erreurs rencontrées et résolution}

\begin{table}[ht]
\centering\small
\caption{Erreurs de la mission 2.}\label{tab:m2-erreurs}
\begin{tabularx}{\linewidth}{@{}>{\raggedright\arraybackslash}p{4.2cm} L L@{}}
\toprule
Symptôme & Cause & Résolution\\ \midrule
Deux pipelines par push sur une branche avec MR ouverte (analyse des règles) & un job sans \code{rules} vaut \code{except: merge_requests} & \code{workflow:rules} et \code{rules: - when: on_success}\\
Archive \code{cpack} sans \fic{BUILDINFO.txt} (\fic{logs/m2-package-cpack.txt}) & le chemin du build est figé dans le cache CMake : il pointe vers le répertoire du job \code{build} & archive construite avec \code{tar}, vérifiée par \code{tar tzf} et \code{sha256sum -c}\\
\code{package} pouvait publier un binaire non testé & \code{needs: [build]} contourne l'ordre des stages & \code{needs: [build, test]}, sans les artefacts de \code{test}\\
Masquage invisible en local & l'émulateur n'implémente pas le masquage & vérification sur GitLab.com\\
\bottomrule
\end{tabularx}
\end{table}

\subsection{Bilan}
\begin{description}[font=\normalfont\bfseries, leftmargin=0pt, itemsep=0.2em]
\item[Ce qui fonctionne :] règles validées pour six types de pipelines, lint qui échoue sur
une vraie erreur, archive vérifiable (contenu, \code{BUILDINFO}, SHA-256), secret transmis par
l'entrée standard.
\item[Ce qui ne fonctionne pas :] \ifgitlab\else masquage et protection non observables hors de
GitLab.com ; \fi le binaire de l'archive dépend de l'Open MPI de Debian et n'est pas autonome.
\item[À refaire autrement :] poser \code{workflow:rules} dès le premier job conditionnel, et
produire les archives de version à partir de l'environnement Spack plutôt que d'une image Debian.
\end{description}

% =============================================================================
\section{Mission 3 : du CI vers Slurm}
% =============================================================================

\subsection{Réalisation}

Le job Slurm \fic{slurm/job_ci.sh} (annexe~B) demande 2 nœuds de 2 tâches (nœuds
\code{e2-standard-2} à 2 vCPU), 5 minutes et la partition \code{debug}. Il charge
l'environnement dans l'ordre Lmod, environnement Spack du dépôt, puis système, et affiche la
source retenue ; sur le cluster du Cours 5, sans Lmod, c'est Spack. Il lance
\code{srun --mpi=pmix --kill-on-bad-exit=1} avec \code{EXPECTED_RANKS=$SLURM_NTASKS} et se termine
par \code{SUCCESS} ou \code{FAILED}.

Le script \fic{slurm/submit_and_poll.sh} (annexe~B) soumet le job par \code{sbatch --parsable},
enregistre son identifiant dans \fic{.slurm_jobid}, puis interroge \code{sacct -X} (repli sur
\code{squeue} tant que la base n'a pas l'entrée) et affiche l'état, la raison d'attente et le
temps écoulé. Il reconnaît tous les états terminaux, y compris ceux qu'oublie la boucle des
diapositives (\code{OUT_OF_MEMORY}, \code{NODE_FAIL}, \code{PREEMPTED}...), archive \code{sacct} et la
sortie du job, et rend 0 (succès), 1 (échec du job) ou 75 (échec d'infrastructure).

Côté GitLab (listing~\ref{lst:m3-yaml}), un gabarit \code{.hpc} porte le tag \code{hpc-login},
attend le job \code{test} sans récupérer son binaire, recompile avec le MPI du cluster et annule
le job Slurm dans \code{after_script}. Deux jobs l'utilisent : \code{test_slurm} (suivi, \code{timeout},
\code{retry}) et \code{test_slurm_wait} (\code{sbatch --wait}).

\lstinputlisting[language=yaml, linerange={140-143,152-158,170-185},
  caption={Délais, annulation et jobs Slurm (\texttt{.gitlab-ci.yml}, extrait des lignes 140 à 185).},
  label=lst:m3-yaml]{.gitlab-ci.yml}

Pour atteindre le cluster, le sujet propose un runner installé sur le nœud de connexion ou une
connexion SSH depuis un runner partagé (tableau~\ref{tab:runner}). Nous retenons le runner sur
le nœud de connexion (figure~\ref{fig:archi}) : il n'ouvre aucune connexion entrante, partage
\fic{/home} avec les nœuds de calcul et ne consomme pas de minutes de runner partagé pendant
l'attente en file. La variante SSH (\code{test_slurm_ssh}, \fic{slurm/ci_remote.sh}) est écrite
mais ne s'active que si les variables SSH existent. Le script \fic{scripts/gcp_evidence.sh}
enregistre le runner (exécuteur shell) quand on lui fournit le jeton créé dans l'interface de
GitLab, où sont fixés le tag \code{hpc-login} et la protection.

\begin{table}[ht]
\centering\small
\caption{Deux façons d'atteindre le cluster depuis GitLab.}\label{tab:runner}
\begin{tabularx}{\linewidth}{@{}>{\raggedright\arraybackslash}p{2.4cm} L L@{}}
\toprule
 & Runner sur le nœud de connexion & SSH depuis un runner partagé\\ \midrule
Réseau & connexion sortante uniquement (le runner interroge GitLab en HTTPS) & IP publique et port 22 ouverts ; les IP des runners partagés ne sont pas fixes\\
Secrets & aucun secret de connexion dans GitLab & clé privée en variable \emph{File} et \emph{Protected}, \code{known_hosts}\\
Fichiers & build sur \fic{/home} (NFS), visible des nœuds & \code{rsync} aller-retour des sources et des logs\\
Attente en file & ne coûte rien côté GitLab & minutes de runner partagé consommées\\
Précautions & runner protégé, utilisateur dédié, aucun calcul sur le nœud de connexion & coupure SSH = job perdu ; \code{after_script} dans un shell séparé\\
\bottomrule
\end{tabularx}
\end{table}

\begin{figure}[ht]
\centering
\begin{tikzpicture}[font=\sffamily\footnotesize, x=1cm, y=1cm]
  \node[boite, fill=oiPourpre!12, text width=2.4cm, align=center, minimum height=2.3cm] (gl) at (0,0)
    {\textbf{GitLab.com}\\[3pt] dépôt \texttt{mpiflow}\\ pipeline\\ artefacts};
  \node[boite, fill=oiCiel!14, text width=5.0cm, align=left] (login) at (6.0,0)
    {\textbf{Nœud de connexion}\\[1pt]
     {\scriptsize runner GitLab : exécuteur shell, tag \texttt{hpc-login}, protégé}\\[4pt]
     {\scriptsize 1. activation de l'environnement Spack\\
      2. compilation avec le MPI du cluster\\
      3. soumission (\texttt{sbatch -{}-parsable})\\
      4. suivi (\texttt{sacct}, \texttt{squeue})\\
      5. artefacts : sortie du job, \texttt{sacct}\\
      en \texttt{after\_script} : \texttt{scancel}}};
  \node[boite, fill=oiVert!10, text width=2.6cm, align=center] (ctl) at (12.35,1.0)
    {\textbf{Slurm}\\ {\scriptsize contrôleur et comptabilité}};
  \node[boite, fill=oiOrange!18, text width=1.1cm, align=center, font=\sffamily\scriptsize] (n1) at (11.5,-1.05)
    {nœud 1\\ rangs 0--1};
  \node[boite, fill=oiOrange!18, text width=1.1cm, align=center, font=\sffamily\scriptsize] (n2) at (13.2,-1.05)
    {nœud 2\\ rangs 2--3};
  \node[boite, fill=black!5, minimum width=15.2cm, align=center] (nfs) at (6.25,-2.95)
    {\texttt{/home} partagé (NFS) : dépôt, build, environnement Spack, sorties des jobs};
  \draw[{Stealth[length=2.2mm]}-{Stealth[length=2.2mm]}, thick] (gl.east) --
    node[above, font=\sffamily\scriptsize] {HTTPS}
    node[below, font=\sffamily\scriptsize] {sortant} (login.west);
  \draw[fleche] (login.east |- ctl.west) --
    node[above, font=\scriptsize, align=center] {\texttt{sbatch}\\[-1pt]\texttt{sacct}} (ctl.west);
  \draw[fleche] (ctl.south) -- node[right, font=\sffamily\scriptsize, pos=0.4] {\texttt{srun}}
    ($(n1.north)!0.5!(n2.north)$);
  \draw[black!40, thick] (login.south) -- (login.south |- nfs.north);
  \draw[black!40, thick] ($(n1.south)!0.5!(n2.south)$) -- ($(n1.south)!0.5!(n2.south)$ |- nfs.north);
  \draw[fleche, dashed, black!50] (gl.north) to[out=50,in=140]
    node[above, font=\sffamily\scriptsize, text=black!60] {variante SSH écartée : runner partagé $\to$ port 22 entrant}
    (login.north);
\end{tikzpicture}
\caption{Architecture retenue pour la mission 3. Le runner orchestre depuis le nœud de
connexion et ne calcule jamais : le calcul a lieu sur les nœuds alloués par Slurm. Le
\texttt{/home} partagé rend le binaire compilé et les sorties du job visibles des deux côtés. En
tirets, la variante SSH écartée.}
\label{fig:archi}
\end{figure}

\subsection{Preuves d'exécution}\label{sec:m3-preuves}

Toutes les exécutions de cette section ont lieu sur le mini-cluster, soumises par un utilisateur
sans privilèges ; la figure~\ref{fig:chrono} les replace sur une même échelle de temps. Le cas
nominal est donné par le listing~\ref{lst:m3-ok} : 4 rangs sur 2 nœuds, \code{ExitCode 0:0},
ligne \code{SUCCESS}. Une première version du script interrogeait \code{sacct} dès l'état terminal ;
la comptabilité étant asynchrone, le nom du job et la liste des nœuds n'étaient pas encore à jour
(listing~\ref{lst:m3-race}). Le script attend désormais quelques secondes avant d'archiver.

\lstinputlisting[caption={Soumission et suivi : succès sur 2 nœuds (\texttt{logs/m3-poll-4-success.txt}).},
  label=lst:m3-ok]{logs/m3-poll-4-success.txt}
\lstinputlisting[caption={\texttt{sacct} interrogé trop tôt, puis 3 s plus tard (\texttt{logs/m3-sacct-race.txt}).},
  label=lst:m3-race]{logs/m3-sacct-race.txt}

\begin{figure}[ht]
\centering
\begin{tikzpicture}[x=0.24cm, y=1cm, font=\sffamily\scriptsize]
  \def\lab{-1.2}
  % axe
  \draw[black!60] (0,0.15) -- (46,0.15);
  \foreach \t in {0,5,...,45} {\draw[black!60] (\t,0.15) -- (\t,0.25); \node[below, text=black!70] at (\t,0.15) {\t};}
  \node[below, text=black!70] at (23,-0.2) {temps depuis la soumission (s)};
  % Job 1 : RUNNING 0 -> 313 s, plantage à 1,5 s
  \node[left, align=right] at (\lab,5.4) {job 1 : un rang plante,\\ sans \texttt{-{}-kill-on-bad-exit}};
  \fill[oiBleu!75] (0,5.25) rectangle (44,5.55);
  \draw[decorate, decoration={zigzag, amplitude=1pt, segment length=3pt}, white, line width=1.4pt] (44.3,5.2) -- (44.3,5.6);
  \fill[oiBleu!75] (44.6,5.25) rectangle (46,5.55);
  \node[circle, fill=white, inner sep=0.6pt, text=oiVermillon, font=\bfseries\small] at (1.5,5.4) {$\times$};
  \node[above, text=oiVermillon] at (1.5,5.55) {segfault};
  \node[above left] at (46,5.55) {TIMEOUT à 313 s};
  % Job 2
  \node[left, align=right] at (\lab,4.5) {job 2 : même échec,\\ avec \texttt{-{}-kill-on-bad-exit=1}};
  \fill[oiBleu!75] (0,4.35) rectangle (1,4.65);
  \node[text=oiVermillon, font=\bfseries\small] at (1.6,4.5) {$\times$};
  \node[right] at (2.2,4.5) {FAILED à 1 s};
  % Job 9
  \node[left, align=right] at (\lab,3.6) {job 9 : file saturée\\ par le job 8 (4 CPU)};
  \fill[oiOrange!45] (0,3.45) rectangle (30,3.75);
  \node at (15,3.6) {PENDING (Resources)};
  \fill[oiBleu!75] (30,3.45) rectangle (31,3.75);
  \fill[oiVert] (31.8,3.6) circle[radius=2.2pt];
  \node[right] at (32.3,3.6) {COMPLETED à 31 s};
  % Job 13
  \node[left, align=right] at (\lab,2.7) {job 13 : \texttt{sbatch -{}-wait},\\ CI coupé à 15 s};
  \fill[oiOrange!45] (0,2.55) rectangle (15,2.85);
  \node at (7.5,2.7) {PENDING};
  \draw[black, thick] (15,2.45) -- (15,2.95);
  \node[right, align=left] at (15.4,2.7) {timeout du CI : le job reste en file\\ \texttt{scancel} $\to$ CANCELLED};
  % Job 14
  \node[left, align=right] at (\lab,1.8) {job 14 : panne\\ de \texttt{node2} à 12 s};
  \fill[oiBleu!75] (0,1.65) rectangle (12,1.95);
  \node[text=oiVermillon, font=\bfseries\small] at (12.6,1.8) {$\times$};
  \fill[oiOrange!45] (13.2,1.65) rectangle (18,1.95);
  \draw[black, thick] (18,1.55) -- (18,2.05);
  \node[right, align=left] at (18.4,1.8) {NODE\_FAIL, remis en file par Slurm\\ (BeginTime), puis \texttt{scancel} à 18 s};
  % légende
  \begin{scope}[shift={(0,-1.05)}]
    \fill[oiOrange!45] (0,-0.1) rectangle (1.6,0.1); \node[right] at (1.8,0) {PENDING};
    \fill[oiBleu!75] (8,-0.1) rectangle (9.6,0.1); \node[right] at (9.8,0) {RUNNING};
    \node[text=oiVermillon, font=\bfseries\small] at (16.5,0) {$\times$}; \node[right] at (17.2,0) {échec};
    \fill[oiVert] (22.5,0) circle[radius=2.2pt]; \node[right] at (23.1,0) {succès};
    \draw[black, thick] (28.5,-0.18) -- (28.5,0.18); \node[right] at (29,0) {annulation ou coupure du CI};
  \end{scope}
\end{tikzpicture}
\caption{Chronologie des jobs Slurm des expériences de la mission 3, d'après les heures de
soumission, de début et de fin enregistrées par \texttt{slurmdbd} (\texttt{logs/m3-chronologie.txt})
et le journal de \texttt{slurmd} pour l'instant du plantage. L'axe du job 1 est interrompu : ses
rangs restants sont restés bloqués jusqu'au wall-time de 5 minutes.}
\label{fig:chrono}
\end{figure}

\paragraph{Un rang plante, les autres attendent.} Notre première soumission sur deux nœuds a mis
plus de cinq minutes à échouer (listing~\ref{lst:m3-timeout}) : deux rangs ont planté dans le
transport mémoire partagée \code{vader}, les deux autres sont restés bloqués dans la collective
jusqu'à ce que Slurm tue le job au wall-time. Avec \code{srun --kill-on-bad-exit=1}, le même
échec est détecté en une seconde (listing~\ref{lst:m3-kill}). Le plantage lui-même vient de
l'émulation : les deux nœuds Slurm sont la même machine, et Open MPI y mélange ses segments de
mémoire partagée. Sur un seul nœud, le test passe, et sur deux nœuds avec
\code{OMPI_MCA_btl=self,tcp} aussi (listing~\ref{lst:m3-ok}) ; ce réglage ne concerne que le
mini-cluster.

\lstinputlisting[caption={Échec lent : TIMEOUT Slurm après 313 s (\texttt{logs/extraits/m3-poll-1-timeout.txt}).},
  label=lst:m3-timeout]{logs/extraits/m3-poll-1-timeout.txt}
\lstinputlisting[lastline=21, caption={Échec rapide avec \texttt{-{}-kill-on-bad-exit=1} (\texttt{logs/m3-poll-2-kill.txt}).},
  label=lst:m3-kill]{logs/m3-poll-2-kill.txt}

\paragraph{File saturée.} Un premier job occupe les 4 CPU du cluster ; le job CI reste
\code{PENDING (Resources)} pendant 30\,s et le suivi l'affiche à chaque tour
(listing~\ref{lst:m3-pending}).

\lstinputlisting[caption={Job CI en attente derrière un job qui occupe tout le cluster (\texttt{logs/m3-pending.txt}).},
  label=lst:m3-pending]{logs/m3-pending.txt}

\paragraph{\texttt{sbatch -{}-wait}.} Le code retour est bien propagé (0 pour un succès, 1 pour un échec,
\fic{logs/extraits/m3-wait-ok-fail.txt}), mais rien ne s'affiche pendant l'attente. Surtout, quand
le CI coupe l'attente (un \code{timeout 15} simule ici le \code{timeout:} de GitLab), le job Slurm
reste en file et s'exécuterait plus tard pour rien (listing~\ref{lst:m3-orphan}).

\lstinputlisting[caption={\texttt{sbatch -{}-wait} coupé par le CI : job orphelin (\texttt{logs/m3-wait-orphan.txt}).},
  label=lst:m3-orphan]{logs/m3-wait-orphan.txt}

\paragraph{Panne de nœud.} Nous passons \code{node2} en \code{DOWN} pendant un job : le script classe
l'échec comme une panne d'infrastructure (code 75). Mais Slurm a déjà remis le job en file
(\code{BeginTime}) : sans \code{scancel}, il repartirait en plus de la relance éventuelle du CI
(listing~\ref{lst:m3-nodefail}).

\lstinputlisting[caption={Panne de nœud : verdict d'infrastructure et remise en file par Slurm (\texttt{logs/m3-nodefail.txt}).},
  label=lst:m3-nodefail]{logs/m3-nodefail.txt}

\paragraph{Base de comptabilité arrêtée.} Pendant une session, la base MariaDB du
mini-cluster s'est arrêtée, et avec elle \code{slurmdbd} : \code{sacct} ne répondait plus. La
première version du script a alors affiché \code{INCONNU} pendant onze minutes sans jamais
conclure, exactement le blocage que seul le \code{timeout:} du CI aurait arrêté. Le script
interroge désormais le contrôleur (\code{scontrol show job}, qui garde les jobs terminés
5 minutes) quand ni \code{sacct} ni \code{squeue} ne répondent, et rend le code 75 après six tours
sans état. Relancé avec la base toujours arrêtée, il conclut normalement
(listing~\ref{lst:m3-dbdown}).

\lstinputlisting[caption={Suivi avec la base de comptabilité arrêtée, avant puis après correction (\texttt{logs/m3-sacct-down.txt}).},
  label=lst:m3-dbdown]{logs/m3-sacct-down.txt}

\paragraph{srun sans PMI.} Sans \code{--mpi=pmix}, Open MPI détecte le lancement direct par
\code{srun} et refuse de démarrer (\fic{logs/extraits/m3-singleton.txt}). Le contrôle
\code{EXPECTED_RANKS} reste utile avec d'autres implémentations MPI, qui peuvent démarrer en
\enquote{singletons} de taille 1 et réussir à tort (\fic{logs/m1-singleton.txt}).

\paragraph{Sur le cluster GCP.}
\ifgcp
Le même script de soumission, lancé depuis le nœud de connexion (\gcpinfo), place cette fois
les rangs sur deux machines distinctes (listings~\ref{lst:gcp-cluster} et~\ref{lst:gcp-slurm}).
\lstinputlisting[caption={Le cluster GCP (\texttt{logs/gcp/cluster.txt}).},
  label=lst:gcp-cluster]{logs/gcp/cluster.txt}
\lstinputlisting[caption={Soumission et suivi sur le cluster GCP (\texttt{logs/gcp/test\_slurm.txt}).},
  label=lst:gcp-slurm]{logs/gcp/test_slurm.txt}
\IfFileExists{logs/gcp/runner.txt}{\lstinputlisting[caption={Enregistrement du runner \texttt{hpc-login} (\texttt{logs/gcp/runner.txt}).}]{logs/gcp/runner.txt}}{}
\else
\afaire{exécution de \texttt{scripts/gcp\_evidence.sh} sur le nœud de connexion : la soumission
sur deux nœuds réels s'insérera ici}
\fi

\subsection{Polling ou \texttt{sbatch -{}-wait} ?}

\begin{table}[ht]
\centering\small
\caption{Suivi par polling et \texttt{sbatch -{}-wait}, d'après nos essais.}\label{tab:polling}
\begin{tabularx}{\linewidth}{@{}>{\raggedright\arraybackslash}p{2.6cm} L L@{}}
\toprule
 & Polling (\texttt{submit\_and\_poll.sh}) & \texttt{sbatch -{}-wait}\\ \midrule
Taille & script d'une soixantaine de lignes, réutilisable & une à trois lignes\\
Pendant l'attente & état, raison et temps écoulé à chaque tour & silence total\\
Code retour & vérifié (\code{ExitCode} et ligne \code{SUCCESS}) et classé : code 1, infrastructure 75 & propagé tel quel, sans la cause\\
Après un timeout du CI & identifiant connu, \code{after_script} annule le job & job orphelin, sauf \code{--parsable > .slurm_jobid}\\
Charge & une requête \code{sacct} par intervalle (30\,s) & aucune requête du script\\
\bottomrule
\end{tabularx}
\end{table}

Le tableau~\ref{tab:polling} compare les deux versions. Nous choisissons le polling pour le
pipeline principal et le nightly : la file est imprévisible, et il faut savoir où en était le job
et distinguer les échecs. \code{sbatch --wait} convient aux jobs courts sur une partition de debug
peu chargée et aux scripts qui enchaînent des soumissions (comme la boucle de passage à l'échelle
des diapositives), à condition d'écrire l'identifiant du job pour pouvoir l'annuler.

\subsection{Questions de cadrage}

\subsubsection{Pourquoi le job CI doit-il porter un timeout: alors que le job Slurm a un wall-time ?}

Parce que les deux horloges ne mesurent pas la même chose. Le wall-time de Slurm ne compte que le
temps passé en \code{RUNNING} ; l'attente en \code{PENDING} n'est bornée par rien (30\,s dans notre
expérience, des heures sur un cluster chargé). La durée du job CI vaut attente, plus exécution,
plus latence du suivi. Le job CI peut aussi se bloquer indépendamment de Slurm : la boucle des
diapositives tourne indéfiniment sur un état qu'elle ne connaît pas (\code{NODE_FAIL},
\code{OUT_OF_MEMORY}...) ou si \code{sacct} ne répond plus, et une connexion SSH peut rester
pendante. Pendant ce temps, le job occupe un emplacement de runner, et des minutes facturées
dans la variante SSH.

Le \code{timeout:} doit donc dépasser le wall-time plus une attente raisonnable, sans quoi il
tuerait un job légitime : nous avons 5\,min de wall-time et 45\,min de \code{timeout}. Cela ne
suffit pas : à l'expiration, GitLab arrête le script, et si \code{after_script} n'a pas le temps de
s'exécuter, le job Slurm reste en file. La documentation impose que \code{RUNNER_SCRIPT_TIMEOUT}
et \code{RUNNER_AFTER_SCRIPT_TIMEOUT} ne dépassent pas ensemble le \code{timeout} du job, d'où
40\,min et 3\,min pour 45\,min (listing~\ref{lst:m3-yaml}). Notre premier essai illustre l'autre
sens : un rang planté, les autres bloqués, et seul le wall-time de Slurm a mis fin au gâchis. Les
deux garde-fous sont complémentaires.

\subsubsection{Le job Slurm reste PENDING 20 minutes : que se passe-t-il côté GitLab ?}

Côté GitLab, le job est \enquote{running} : le runner exécute la boucle de suivi. Il occupe un
emplacement du runner, si bien que les autres pipelines attendent, consomme des minutes de
runner partagé dans la variante SSH, et bloque la suite du pipeline (et une MR si le job est
requis). Avec \code{sbatch --wait}, le log reste vide pendant 20 minutes : impossible de distinguer
une file pleine d'un blocage. Si l'attente et l'exécution dépassent le \code{timeout}, le job CI
échoue et le job Slurm reste en file s'il n'est pas annulé (figure~\ref{fig:chrono}, job 13).

Pour le rendre observable, notre script affiche à chaque tour l'heure, le temps écoulé, l'état et
la raison donnée par Slurm (\code{Resources}, \code{Priority}, \code{QOSMaxJobsPerUserLimit},
\code{ReqNodeNotAvail}...), puis une fois l'heure de démarrage estimée par \code{squeue --start},
qui reste \code{N/A} tant que l'ordonnanceur n'a pas calculé de prévision. L'identifiant est
conservé, et \code{sacct} et la sortie du job sont archivés même en cas d'échec. On peut aller plus
loin avec \code{sbatch --test-only} pour estimer le démarrage avant de soumettre, une partition ou
une QoS réservée au CI, et une alerte au-delà d'un seuil d'attente. Sur un cluster
cloud dont les nœuds sont créés à la demande, une partie de l'attente vient en outre du
démarrage des VM, que \code{sinfo} signale par un suffixe sur l'état des nœuds.

\subsubsection{Le nightly échoue à 2 h du matin : que se demander avant de relancer ?}

Trois questions, dans l'ordre :
\begin{enumerate}
\item \textbf{Quel type d'échec ?} La raison donnée par GitLab (\code{script_failure} contre
\code{runner_system_failure}, \code{runner_interrupted} ou timeout) et l'état Slurm (\code{FAILED} avec
un \code{ExitCode}, contre \code{NODE_FAIL}, \code{BOOT_FAIL}, \code{PREEMPTED} ; \code{TIMEOUT} de Slurm
contre timeout de GitLab). Notre script y répond directement : code 1 ou code 75.
\item \textbf{Ai-je les logs ?} Les artefacts conservés même en cas d'échec (\code{when: always}) :
\fic{ci-JOBID.out}, \fic{sacct-JOBID.txt} et le rapport de tests, puis \code{sacct -j JOBID
--format=ALL} pour la mémoire et le nœud. Sans logs, c'est l'observabilité qu'il faut réparer
avant de relancer.
\item \textbf{Qu'est-ce qui a changé depuis le dernier nightly réussi ?} Commits fusionnés sur
\code{main}, différences du \fic{spack.lock}, maintenance du cluster, quota ou crédits GCP épuisés.
Si rien n'a changé dans le dépôt, l'infrastructure est la première suspecte.
\end{enumerate}
On peut en ajouter une quatrième : quelqu'un doit-il vraiment se lever ? Un nightly en échec
bloque rarement une version ; les alertes immédiates se réservent à \code{main} et aux tags.

Un échec de code est déterministe et lié à un commit : test en échec, \code{FAILED} avec un
\code{ExitCode} non nul, assertion, segfault dans notre code. Le relancer ne sert à rien, masque les
tests instables et consomme des heures de calcul : il faut corriger. Un échec d'infrastructure est
transitoire et indépendant du code : runner perdu, nœud défaillant, VM impossible à créer (quota
ou ressources de la zone épuisés), préemption, NFS. Une relance automatique, limitée à une, est
alors légitime. D'où notre configuration : \code{retry} avec \code{max: 1},
\code{when: [runner_system_failure, runner_interrupted]} et \code{exit_codes: 75}. La documentation
récente indique que GitLab 19.1 a séparé de \code{runner_system_failure} les cas
\code{runner_interrupted} et \code{runner_external_dependency_failure} : ne lister que l'ancienne
valeur ne couvre plus un runner redémarré. Enfin, Slurm remet lui-même en file un job après
\code{NODE_FAIL} : il faut l'annuler avant que le CI ne relance, sinon le test tourne deux fois.

\subsection{Erreurs rencontrées et résolution}

\begin{table}[ht]
\centering\small
\caption{Erreurs de la mission 3.}\label{tab:m3-erreurs}
\begin{tabularx}{\linewidth}{@{}>{\raggedright\arraybackslash}p{4.2cm} L L@{}}
\toprule
Symptôme & Cause & Résolution\\ \midrule
\code{Could not establish a configuration source} & Slurm installé sans \fic{slurm.conf} & écriture de \fic{slurm.conf} et \fic{slurmdbd.conf}, démarrage des démons\\
Segfault dans \code{mca_btl_vader} sur 2 nœuds & émulation : deux nœuds Slurm sur une même machine & 1 nœud : OK ; 2 nœuds : \code{OMPI_MCA_btl=self,tcp} (mini-cluster seulement)\\
Rangs survivants bloqués 5\,min & \code{srun} laisse tourner les autres tâches & \code{--kill-on-bad-exit=1} : échec en 1\,s\\
\code{OMPI was not built with SLURM's PMI support} & \code{srun} sans plugin PMI & \code{srun --mpi=pmix}\\
\code{sacct} : \code{JobName=allocation}, nœuds incomplets & enregistrement asynchrone dans \code{slurmdbd} & attente de 5\,s avant l'archivage\\
Suivi bloqué sur \code{INCONNU} pendant 11\,min (listing~\ref{lst:m3-dbdown}) & base MariaDB, donc \code{slurmdbd}, arrêtée & repli sur \code{scontrol show job}, code 75 après six tours sans état\\
Job relancé par Slurm après \code{NODE_FAIL} & remise en file automatique par défaut & \code{scancel} en \code{after_script}\\
\code{after_script} non exécuté après un timeout (documentation) & le timeout global coupe aussi \code{after_script} & 40\,min + 3\,min < \code{timeout} de 45\,min\\
Jobs \code{tags: [hpc-login]} en attente sans fin sur GitLab.com (anticipé) & aucun runner ne porte ce tag & jobs créés seulement si \code{HPC_ENABLED} vaut \code{true}\\
\bottomrule
\end{tabularx}
\end{table}

\subsection{Bilan}
\begin{description}[font=\normalfont\bfseries, leftmargin=0pt, itemsep=0.2em]
\item[Ce qui fonctionne :] soumission, suivi, verdict classé (0, 1 ou 75) et annulation des
jobs orphelins, testés sur un vrai Slurm avec comptabilité, file saturée et panne de nœud.
\item[Ce qui ne fonctionne pas :] \ifgcp si la base de comptabilité reste arrêtée plus de 5
minutes après la fin du job, le contrôleur l'a oublié et le code retour est perdu : le script
conclut alors à un échec. \else pas encore
d'exécution sur des nœuds réellement distincts ; le transport mémoire partagée n'a pas pu être
testé en conditions réelles.\fi
\item[À refaire autrement :] tester les cas d'échec (plantage, file pleine, panne) avant le cas
nominal : chacun a révélé un défaut du script.
\end{description}

% =============================================================================
\section{Mission 4 : Spack, l'environnement comme code}
% =============================================================================

\subsection{Réalisation}

Spack est cloné dans le répertoire personnel au tag \code{v1.2.2} plutôt que sur la branche
\code{develop}, car la version de Spack fait partie de ce qu'il faut figer. \code{spack compiler
find} détecte le GCC 13.3 du système, déclaré comme paquet externe : Spack ne le recompile pas.
L'environnement \fic{spack/} demande \code{openmpi@5 schedulers=slurm %gcc} et \code{cmake %gcc}
(listing~\ref{lst:m4-yaml}). Deux réglages se sont imposés à l'expérience. D'une part, une cible
générique \code{x86_64_v3} : \code{spack arch} donne \code{icelake} sur notre machine, alors que les
VM \code{e2} reçoivent, selon la disponibilité, un Xeon Haswell, Broadwell ou Skylake ou un AMD EPYC
Rome, dont \code{x86_64_v3} est le dénominateur commun. D'autre part, le Slurm du système est
déclaré externe et non constructible : sans cela, \code{schedulers=slurm} fait compiler par Spack
un Slurm complet, avec Python et GLib (76 paquets à construire au lieu de 46).

\lstinputlisting[language=yaml, style=numerote,
  caption={Manifeste de l'environnement (\texttt{spack/spack.yaml}).},
  label=lst:m4-yaml]{spack/spack.yaml}

La concrétisation produit 49 specs, dont 46 à construire et 3 externes : GCC, la glibc et Slurm
(listing~\ref{lst:m4-conc}). Les deux fichiers sont commités dans \fic{spack/} ; la vue et l'état
local (\fic{spack/.spack-env/}) sont ignorés par Git.

\lstinputlisting[lastline=17, caption={Début de la concrétisation (\texttt{logs/m4-concretize.txt}) :
\texttt{-} à construire, \texttt{[e]} externe.},
  label=lst:m4-conc]{logs/m4-concretize.txt}

\ifgcp
Sur le nœud de connexion GCP, le script \fic{scripts/gcp_evidence.sh} adapte le Slurm externe
au cluster, concrétise de nouveau l'environnement (le verrou commité est celui-là) et l'installe
depuis les sources en \spackinstalltime{} (listing~\ref{lst:gcp-spack}) : c'est le coût que
chaque pipeline paierait sans buildcache.
\lstinputlisting[caption={Concrétisation et installation sur le cluster GCP (\texttt{logs/gcp/spack.txt}, \texttt{spack-install.txt}).},
  label=lst:gcp-spack]{logs/gcp/spack.txt}
\lstinputlisting{logs/gcp/spack-install.txt}
\else
\afaire{exécution de \texttt{scripts/gcp\_evidence.sh} : concrétisation sur le cluster cible et
durée de l'installation}
\fi

\subsection{spack.yaml et spack.lock}

\begin{table}[ht]
\centering\small
\caption{Ce que décrivent le manifeste et le verrou.}\label{tab:yaml-lock}
\begin{tabularx}{\linewidth}{@{}>{\raggedright\arraybackslash}p{2.4cm} L L@{}}
\toprule
 & \texttt{spack.yaml} (manifeste) & \texttt{spack.lock} (verrou)\\ \midrule
Contenu & des intentions : specs abstraites, contraintes, configuration (cible, externes, vue) & le résultat de la concrétisation : le graphe complet\\
Versions & une contrainte sur OpenMPI, rien sur les dépendances & toutes : \code{openmpi@5.0.10}, \code{pmix@6.1.0}, \code{cmake@3.31.11}...\\
Variants & ceux qu'on impose (\code{schedulers=slurm}) & tous, y compris les valeurs par défaut\\
Compilateurs & \texttt{\%gcc} & \code{gcc@13.3.0}, nœud du graphe (dépendance de compilation)\\
Identité & -- & hash de chaque nœud, arêtes typées (\code{build}, \code{link}), OS et cible\\
Forme & 27 lignes de YAML commenté, écrites à la main & JSON généré : 64\,Ko, 49 specs\\
\bottomrule
\end{tabularx}
\end{table}

\begin{figure}[ht]
\centering
\begin{tikzpicture}[font=\sffamily\scriptsize, x=1cm, y=1cm,
    spec/.style={draw=oiBleu, fill=oiBleu!10, rounded corners=2pt, align=center,
      inner sep=3pt, text width=1.5cm},
    ext/.style={spec, draw=black!45, fill=black!6, dashed},
    racine/.style={draw=oiBleu, fill=oiBleu!28, rounded corners=2pt, align=center,
      inner sep=4pt, font=\sffamily\footnotesize}]
  \newcommand{\spk}[3]{\textbf{#1}\\ @#2\\ {\ttfamily #3}}
  \node[racine] (ompi) at (4.8,2.0) {\textbf{openmpi@5.0.10}\\ {\ttfamily iftmwb3}};
  \node[racine] (cmake) at (12.6,2.0) {\textbf{cmake@3.31.11}\\ {\ttfamily nxsxsy2}};
  \node[spec] (pmix)  at (0.9,0)   {\spk{pmix}{6.1.0}{2af34qh}};
  \node[spec] (prrte) at (2.85,0)  {\spk{prrte}{4.1.0}{s5tkl45}};
  \node[spec] (hwloc) at (4.8,0)   {\spk{hwloc}{2.13.0}{3hmup55}};
  \node[spec] (lev)   at (6.75,0)  {\spk{libevent}{2.1.12}{j4kmbgk}};
  \node[ext]  (slurm) at (8.7,0)   {\spk{slurm}{23.11.4}{externe /usr}};
  \node[ext]  (gcc)   at (10.65,0) {\spk{gcc}{13.3.0}{externe /usr}};
  \node[spec] (curl)  at (12.6,0)  {\spk{curl}{8.20.0}{i6kir7u}};
  \node[spec] (ncur)  at (14.55,0) {\spk{ncurses}{6.6}{h5ybcvi}};
  \foreach \n in {pmix,prrte,hwloc,lev,slurm}
    \draw[-{Stealth[length=1.6mm]}, black!70] (ompi.south) -- (\n.north);
  \draw[-{Stealth[length=1.6mm]}, black!70, dashed] (ompi.south) -- (gcc.north);
  \foreach \n in {curl,ncur} \draw[-{Stealth[length=1.6mm]}, black!70] (cmake.south) -- (\n.north);
  \draw[-{Stealth[length=1.6mm]}, black!70, dashed] (cmake.south) -- (gcc.north);
  \node[text=black!65] at (7.7,-1.05) {+ 39 autres specs (autotools, perl, openssh, zlib-ng, numactl...)};
  \begin{scope}[shift={(0.1,-1.65)}]
    \draw[-{Stealth[length=1.6mm]}, black!70] (0,0) -- (0.7,0) node[right, black] {dépendance de build et d'édition de liens};
    \draw[-{Stealth[length=1.6mm]}, black!70, dashed] (6.5,0) -- (7.2,0) node[right, black] {build seulement};
    \node[draw=black!45, fill=black!6, dashed, rounded corners=1pt, minimum width=0.6cm, minimum height=0.25cm, inner sep=0pt] at (10.9,0) {};
    \node[right] at (11.25,0) {externe, non reconstruit};
  \end{scope}
\end{tikzpicture}
\caption{Extrait du graphe enregistré dans \texttt{spack.lock} : les deux specs racines et
quelques dépendances directes, avec les sept premiers caractères de leur hash. Les externes
(GCC, Slurm) sont propres à la machine qui a concrétisé.}
\label{fig:lock}
\end{figure}

Le tableau~\ref{tab:yaml-lock} et la figure~\ref{fig:lock} résument ce que contient chaque
fichier ; le listing~\ref{lst:m4-lock} montre la spec concrète d'OpenMPI telle qu'elle est
enregistrée.

\lstinputlisting[caption={Extrait de \texttt{spack.lock}, mis en forme (\texttt{logs/m4-lock-excerpt.txt}).},
  label=lst:m4-lock]{logs/m4-lock-excerpt.txt}

\subsection{Questions de cadrage}

\subsubsection{Syntaxe package@version\%compiler+variant\textasciitilde variant et spec de HDF5}

\code{@} fixe une version ou un intervalle (\code{@1.14.6}, \code{@1.14:}), \code{%} le compilateur,
\code{+} active un variant booléen et \code{~} le désactive ; \code{clé=valeur} règle un variant à
valeurs (\code{schedulers=slurm}, \code{target=x86_64_v3}) et \code{^} contraint une dépendance
(\code{^openmpi@5}). Depuis Spack~1.0, le compilateur est une dépendance comme une autre : la
sortie affiche \code{%c,cxx=gcc@13.3.0}. La spec demandée s'écrit \code{hdf5+mpi+fortran~cxx %gcc},
ou, entièrement figée, \code{hdf5@1.14.6+mpi+fortran~cxx %gcc@13.3.0 ^openmpi@5.0.10}. On vérifie
avec \code{spack spec} qu'elle est bien comprise (listing~\ref{lst:m4-hdf5}). D'après \code{spack info
hdf5}, \code{cxx} est déjà désactivé par défaut et \code{mpi} activé : les écrire quand même protège
la spec contre un changement de valeur par défaut dans une future version de la recette. On
remarque aussi la cible \code{icelake} hors de l'environnement, contre \code{x86_64_v3} dans
l'environnement.

\lstinputlisting[caption={Vérification de la spec par \texttt{spack spec} (\texttt{logs/m4-hdf5.txt}).},
  label=lst:m4-hdf5]{logs/m4-hdf5.txt}

\subsubsection{Pourquoi spack.lock suffit-il à reproduire l'environnement, et pas spack.yaml seul ?}\label{sec:lock}

\fic{spack.yaml} décrit des intentions, et le résultat de leur concrétisation dépend du
contexte : version de Spack et des recettes (les versions préférées évoluent : aujourd'hui
OpenMPI 5.0.10 et PMIx 6.1.0, autre chose dans six mois), compilateurs disponibles, paquets
externes déclarés, micro-architecture de la machine (\code{icelake} ici si l'on ne force pas la
cible), réglages du solveur. Deux machines, ou deux dates, donnent deux environnements différents.
\fic{spack.lock} fige au contraire le graphe complet : versions, variants, compilateurs, cible,
arêtes et hash. \code{spack env create mpiflow spack/spack.lock} suivi de \code{spack install}
réinstalle exactement ce graphe sans relancer le solveur, et le hash identifie la configuration :
même hash, même build, même entrée de buildcache.

Le verrou a toutefois des limites, que l'exercice fait apparaître. Il suppose que la version de
Spack sait encore relire ses recettes (d'où le tag \code{v1.2.2}) et que les sources restent
téléchargeables (leurs sommes de contrôle sont vérifiées). Surtout, il contient des éléments
propres à la machine : l'OS (\code{ubuntu24.04}) et les externes (GCC et Slurm 23.11.4 dans
\fic{/usr}, figure~\ref{fig:lock}). Notre verrou local démontre le mécanisme, mais le verrou de
référence doit être produit sur le cluster cible, dont le nœud de connexion et les nœuds de calcul
partagent la même image, puis commité. On en déduit qu'il est reproductible \enquote{à machine
identique}, pas sur n'importe quel poste.

\subsubsection{Quel est l'intérêt du buildcache pour un pipeline qui installe l'environnement à chaque exécution ?}

Un job CI part d'un environnement vierge : \code{spack install} y recompilerait les 46 paquets de
notre environnement (OpenMPI, PMIx, PRRTE, hwloc, CMake, autotools...) à chaque pipeline, soit
des dizaines de minutes à quelques heures sur 2 vCPU, bien au-delà de notre \code{timeout} de
45\,min. Un buildcache stocke les binaires indexés par leur hash : l'installation devient un
téléchargement suivi d'une relocalisation, en quelques minutes. Comme les hash du verrou sont
stables, le cache est atteint exactement, et seuls les nœuds modifiés se recompilent. Le pipeline
de \code{main} peut publier les nouveaux binaires (\code{spack buildcache push}), signés par la clé GPG
du projet (variable \code{SPACK_GPG_KEY}, de type \emph{File} et protégée), dans un registre OCI (le
registre de conteneurs du projet GitLab) ou dans un bucket GCS. Le cluster et le CI installent
alors les mêmes binaires. Par défaut, Spack ne connaît qu'un miroir de sources
(\code{spack-public [s]}) : un buildcache se configure explicitement.

\subsection{Erreurs rencontrées et résolution}

\begin{table}[ht]
\centering\small
\caption{Erreurs de la mission 4.}\label{tab:m4-erreurs}
\begin{tabularx}{\linewidth}{@{}>{\raggedright\arraybackslash}p{4.2cm} L L@{}}
\toprule
Symptôme & Cause & Résolution\\ \midrule
Échec de l'amorçage de \code{clingo} (\code{ghcr.io} injoignable depuis notre réseau) & Spack télécharge son solveur au premier \code{spack spec} & \code{pip install clingo} : Spack utilise un clingo importable par son Python\\
76 paquets à construire, dont \code{slurm@25-05-1-1}, Python, GLib & \code{schedulers=slurm} ajoute Slurm comme dépendance & Slurm du système déclaré externe et non constructible : 46 paquets\\
\code{spack external find slurm} : \code{Could not establish a configuration source} & la détection lance \code{salloc --version}, qui échoue sans \fic{slurm.conf} & détection relancée une fois le mini-cluster configuré\\
Commentaires de \fic{spack.yaml} effacés & \code{spack external find} réécrit le manifeste & manifeste réécrit à la main\\
Cible \code{icelake} détectée & concrétisation pour le processeur de la machine & \code{packages: all: target: [x86_64_v3]}\\
\code{spack install} impossible en local & pas d'accès aux miroirs de sources & installation sur le nœud de connexion GCP (\spackinstalltime)\\
\bottomrule
\end{tabularx}
\end{table}

\subsection{Bilan}
\begin{description}[font=\normalfont\bfseries, leftmargin=0pt, itemsep=0.2em]
\item[Ce qui fonctionne :] environnement concrétisé (49 specs) avec une cible portable et Slurm
externe ; \fic{spack.yaml} et \fic{spack.lock} commités ; spec de HDF5 vérifiée.
\item[Ce qui ne fonctionne pas :] \ifgcp pas de buildcache : chaque installation repart des
sources (\spackinstalltime). \else installation, et build du test avec l'OpenMPI de Spack, pas
encore faits (sur GCP).\fi
\item[À refaire autrement :] concrétiser directement sur le nœud de connexion, et mettre en place
le buildcache avant d'appeler \code{spack install} depuis le CI.
\end{description}

% =============================================================================
\section{Conclusion et auto-diagnostic}
% =============================================================================

La séance répond à la question qui la sous-tend : un pipeline HPC ne se contente pas de
transposer le CI du cloud. Les tests rapides tournent sur des runners partagés, mais ce qu'ils
valident s'arrête à la logique du code ; l'intégration avec le MPI, le lanceur et le réseau du
cluster passe par Slurm, avec deux horloges à gérer, des jobs orphelins à annuler et des échecs
à classer avant de relancer. Spack apporte la pièce manquante, un environnement versionné à côté
du code, à condition de le concrétiser sur la machine cible.

Plusieurs réponses aux questions de cadrage sont sorties nuancées de l'expérience : oublier
\code{needs:} ne casse pas un pipeline GitLab, le test MPI des diapositives passe toujours en
\code{Release}, et un rang qui plante ne fait pas échouer un job Slurm avant le wall-time. Les
principales difficultés ont été techniques : les deux nœuds émulés sur une même machine ont fait
planter le transport mémoire partagée d'Open MPI, et la comptabilité asynchrone de Slurm renvoyait
des informations incomplètes juste après la fin d'un job.

\subsection{Acquis}
\begin{itemize}
\item La sémantique réelle d'un pipeline : stages et DAG (\code{needs}), transmission des
artefacts, types de pipelines et \code{workflow:rules}, conditions avec \code{rules}.
\item Les trois protections d'un secret (masqué, protégé, type \emph{File}), ce qu'aucune ne
garantit seule, et le passage par l'entrée standard.
\item L'orchestration du CI vers Slurm : \code{timeout} et wall-time, états terminaux, annulation
des jobs orphelins, distinction entre échec de code et échec d'infrastructure.
\item Le manifeste et le verrou de Spack, les paquets externes et la cible micro-architecture.
\item Une méthode : provoquer volontairement chaque échec avant de faire confiance au cas nominal.
\end{itemize}

\subsection{Points à retravailler}
\begin{itemize}
\item L'enregistrement et la sécurisation d'un runner sur le cluster GCP (utilisateur dédié,
runner protégé, \fic{/home} partagé).
\item Le buildcache de Spack en pratique : publication, signature GPG, utilisation depuis le CI.
\item La génération de modules Lmod depuis Spack (\code{spack module lmod refresh}).
\item Le runner auto-hébergé de GitHub Actions et les conteneurs Apptainer (cours suivant).
\end{itemize}

\subsection{Utilisation dans un projet d'imagerie médicale}

Prenons un code de segmentation 3D (ITK, MPI ou GPU) appliqué à des volumes IRM ou scanner.
\begin{itemize}
\item \textbf{À chaque push :} lint, tests unitaires et un test de bout en bout sur un petit
fantôme synthétique ou un volume anonymisé de quelques mégaoctets, dans un conteneur.
\item \textbf{Chaque nuit, via Slurm :} segmentation d'un jeu de référence ; le test compare le
coefficient de Dice à une segmentation de référence avec une tolérance, plutôt qu'une égalité bit
à bit, car les réductions MPI et GPU ne sont pas déterministes au bit près.
\item \textbf{Traçabilité :} chaque résultat archivé avec le SHA du commit, le hash du
\fic{spack.lock} (ITK, VTK, HDF5, PyTorch), le JobID Slurm et les sommes SHA-256 des volumes
produits, ce qui permet de rejouer une figure de publication.
\item \textbf{Données de santé :} aucune donnée de patient dans les artefacts, les logs ou les
caches du CI (RGPD) ; seulement des données anonymisées ou synthétiques, un accès au stockage par
variable protégée, et les jobs réels uniquement sur le cluster habilité.
\end{itemize}

% =============================================================================
\clearpage
\section*{Références}
% =============================================================================
\begin{raggedright}
\begin{enumerate}[label=\arabic*.]
\item L. A. Steffenel, \emph{Virtualisation et cloud pour le HPC, Cours 6 : DevOps}, M2 CHPS, Université de Reims Champagne-Ardenne, 2026.
\item GitLab, \emph{CI/CD YAML syntax reference} (\texttt{needs}, \texttt{retry}, \texttt{after\_script}, \texttt{timeout}), \url{https://docs.gitlab.com/ci/yaml/}.
\item GitLab, \emph{Job artifacts}, \url{https://docs.gitlab.com/ci/jobs/job_artifacts.html}.
\item GitLab, \emph{Specify when jobs run with rules}, \url{https://docs.gitlab.com/ci/jobs/job_rules/}.
\item GitLab, \emph{Deprecated keywords} (\texttt{only}, \texttt{except}), \url{https://docs.gitlab.com/ci/yaml/deprecated_keywords/}.
\item GitLab, \emph{CI/CD variables}, \url{https://docs.gitlab.com/ci/variables/}.
\item GitLab, \emph{Configuring runners}, \url{https://docs.gitlab.com/ci/runners/configure_runners/}.
\item Spack, documentation de la version 1.x, \url{https://spack.readthedocs.io/}.
\item Google Cloud, annonce des machines E2, \url{https://cloud.google.com/blog/products/compute/google-compute-engine-gets-new-e2-vm-machine-types}.
\item GitHub, \texttt{actions/upload-artifact}, \url{https://github.com/actions/upload-artifact}.
\item gitlab-ci-local, \url{https://github.com/firecow/gitlab-ci-local}.
\end{enumerate}
\end{raggedright}

% =============================================================================
\clearpage
\section*{Annexe A : pipeline \emph{.gitlab-ci.yml}}
% =============================================================================

Ce fichier définit le pipeline complet de la figure~\ref{fig:pipeline} : jobs des missions 1 et
2 pour les runners partagés, et jobs Slurm de la mission 3 pour le runner du nœud de connexion.

\lstinputlisting[language=yaml, style=numerote]{.gitlab-ci.yml}

\clearpage
\section*{Annexe B : scripts Slurm \emph{job\_ci.sh} et \emph{submit\_and\_poll.sh}}

Le premier script est le job soumis à Slurm ; le second le soumet, suit son état et rend un
verdict exploitable par GitLab. Ils s'exécutent depuis la racine du dépôt.

\lstinputlisting[language=bash, style=numerote, title={\small\texttt{slurm/job\_ci.sh}}]{slurm/job_ci.sh}
\lstinputlisting[language=bash, style=numerote, title={\small\texttt{slurm/submit\_and\_poll.sh}}]{slurm/submit_and_poll.sh}

\clearpage
\section*{Annexe C : test MPI \emph{test\_mpi.c}}

Le test est compilé par CMake et lancé par CTest (à 1 et 4 rangs), par \code{mpirun} dans le
conteneur CI et par \code{srun} sur le cluster.

\lstinputlisting[language=C, style=numerote]{tests/test_mpi.c}

\end{document}
